\section{Background and Problem Formulation}

\subsection{Finite-sample scaling spectra}

Let \(x_t\), \(t=1,\ldots,T\), denote a finite sequence of observed
increments. In the synthetic experiments below, \(x_t\) is a standardized
return or increment sequence rather than an already integrated path. A common
route to multifractal analysis is to examine how empirical moments vary across
aggregation scales. For a scale \(a\), define the aggregated increment
\begin{equation}
  \Delta_a X_t=\sum_{i=0}^{a-1}x_{t+i}.
\end{equation}
A structure-function estimate is then based on
\begin{equation}
  S_q(a)=\frac{1}{T-a+1}\sum_{t=1}^{T-a+1}
  |\Delta_a X_t|^q ,
  \label{eq:structure-function}
\end{equation}
and, over a finite scale set \(\mathcal A\), the scaling exponent is estimated
from the approximate relation
\begin{equation}
  S_q(a)\approx C_q a^{\zeta(q)}.
  \label{eq:scaling-law}
\end{equation}
In practice, \(\zeta(q)\) is obtained by regressing \(\log S_q(a)\) on
\(\log a\) for each \(q\). The deterministic estimators used in the present
study include ordinary least-squares regression, trimmed regression over
scales, bootstrap resampling of scale subsets, and a lightly smoothed
\(\zeta(q)\) curve. These choices are deliberately simple: they allow us to
ask whether the curvature signal is identifiable before attributing failure to
a particular neural architecture.

The finite-sample setting is crucial. The scale set is short, the \(q\)-grid is
finite, and high-order moments are sensitive to rare events. In the main
experiments the default grid is
\[
q\in\{0.5,1.0,1.5,2.0,2.5,3.0\},
\]
and the default scale set is drawn from powers of two such as
\(\{2,4,8,16,32,64\}\), subject to the sample length. These choices are not
claimed to be universal. They are the tested regime in which CMIN-SR is
evaluated.

\subsection{From scaling exponents to singularity spectra}

The scaling function \(\zeta(q)\) is one representation of a multifractal
spectrum. In the classical multifractal formalism, a singularity spectrum can
also be obtained through a Legendre-type transformation,
\begin{equation}
  \alpha(q)=\frac{d\zeta(q)}{dq}, \qquad
  f(\alpha(q))=q\alpha(q)-\zeta(q)+1,
  \label{eq:legendre-spectrum}
\end{equation}
for one-dimensional signals under the usual formal assumptions
\cite{halsey1986fractal,chhabra1989direct}. This paper does not claim a
validated empirical reconstruction of \(f(\alpha)\) for arbitrary real data.
Instead, \(f(\alpha)\) is used to connect the analytic MRW projection family
with standard multifractal geometry. The reported empirical diagnostics focus
on \(\zeta(q)\), projection residuals, curvature, and stability.

\subsection{Monofractal and MRW projection families}

A monofractal spectrum is represented by the linear family
\begin{equation}
  \zeta_{\mathrm{mono}}(q)=qH,
  \label{eq:mono-zeta}
\end{equation}
where \(H\) is a scaling exponent. This family is an essential competitor:
stable scaling by itself is not evidence of multifractality. Fractional
Gaussian noise, for example, can exhibit stable monofractal scaling while
having no nonlinear curvature in \(\zeta(q)\)
\cite{mandelbrot1968fractional}.

The MRW reference family used here follows the parabolic convention
\begin{equation}
  \zeta_{\mathrm{MRW}}(q)=qH-\frac{1}{2}\lambda^2 q(q-2),
  \label{eq:mrw-zeta}
\end{equation}
where \(H\) controls the leading linear component and \(\lambda^2\) controls
spectral curvature \cite{bacry2001multifractal,bacry2003log}. The phrase
``MRW-compatible'' in this paper always refers to compatibility with this
projection family under the selected \(q\)-grid and residual criterion. The
projected coordinate \(\lambda^2_{\mathrm{proj}}\) is not evidence that the
observed time series was generated by an MRW mechanism. Other MRW
parameterizations exist; all reported curvature coordinates in this paper use
the convention in Eq.~\eqref{eq:mrw-zeta}.

\subsection{Projection residuals as diagnostic quantities}

Given an empirical spectrum \(\widehat{\zeta}(q)\), CMIN-SR computes two
least-squares projection coordinates: one for the monofractal family and one
for the MRW family. In schematic form,
\begin{equation}
  \widehat H_{\mathrm{mono}}
  =\arg\min_H
  \sum_{q\in\mathcal Q}
  \left[\widehat{\zeta}(q)-qH\right]^2,
  \label{eq:mono-proj}
\end{equation}
and
\begin{equation}
  (\widehat H_{\mathrm{proj}},\widehat{\lambda}^2_{\mathrm{proj}})
  =\arg\min_{H,\lambda^2}
  \sum_{q\in\mathcal Q}
  \left[
  \widehat{\zeta}(q)
  -
  qH+\frac{1}{2}\lambda^2q(q-2)
  \right]^2 .
  \label{eq:mrw-proj}
\end{equation}
The implementation clamps the MRW curvature coordinate to a nonnegative
admissible range. The corresponding residual norms are used together with
curvature diagnostics. A low MRW residual is not sufficient by itself; the
monofractal residual must also be considered. If the two residuals are nearly
equal and \(\lambda^2_{\mathrm{proj}}\) is small, the sample is better
described as a boundary case rather than as a strong MRW case.
Here, ``boundary'' is an operational diagnostic label for spectra that lie
near the monofractal/MRW projection boundary in the tested finite-sample
geometry; it is not introduced as a separate theoretical process class.

\subsection{The diagnostic inverse problem}

The finite-sample inverse problem studied here is therefore
\begin{equation}
x(t)\longrightarrow
\widehat{\zeta}(q)\longrightarrow
\{\widehat H_{\mathrm{mono}},
  \widehat H_{\mathrm{proj}},
  \widehat{\lambda}^2_{\mathrm{proj}},
  r_{\mathrm{mono}},
  r_{\mathrm{MRW}}\}
\longrightarrow \text{diagnostic scores}.
\end{equation}
The purpose is not to recover a latent parameter at all costs. The purpose is
to decide whether the finite empirical spectrum gives stable evidence for a
linear monofractal explanation, a low-curvature boundary explanation, a curved
MRW-family explanation, or an unstable/tail-driven pattern.

\begin{figure}[t]
  \centering
  \includegraphics[width=0.98\linewidth]{fig1_journal_cmin_sr_pipeline.pdf}
  \caption{CMIN-SR diagnostic pipeline. Raw finite signals are mapped to
  empirical spectra, projected onto monofractal and MRW spectral families, and
  interpreted through residual geometry, curvature, boundary, and instability
  diagnostics. The output is organized diagnostic evidence, not proof of an
  MRW data-generating mechanism.}
  \label{fig:framework}
\end{figure}
